\documentclass[11pt]{article}

\usepackage[margin=1in]{geometry}
\usepackage[T1]{fontenc}
\usepackage[utf8]{inputenc}
\usepackage{lmodern}
\usepackage{microtype}
\usepackage{booktabs}
\usepackage{array}
\usepackage{longtable}
\usepackage{enumitem}
\usepackage{xcolor}
\usepackage{hyperref}
\usepackage{amsmath}
\usepackage{amssymb}

\hypersetup{
  colorlinks=true,
  linkcolor=blue!55!black,
  citecolor=blue!55!black,
  urlcolor=blue!55!black
}

\setlist[itemize]{topsep=3pt,itemsep=2pt,leftmargin=1.4em}
\setlist[enumerate]{topsep=3pt,itemsep=2pt,leftmargin=1.7em}

\newcommand{\code}[1]{\texttt{\detokenize{#1}}}
\newcommand{\risk}{\operatorname{risk}}

\title{\textbf{Methodology Memo}\\
\large Behavioural Risk Intelligence for Agentic Payments}
\author{Bloomsbury Tech -- Agentic Payments Hackathon Build}
\date{April 25, 2026}

\begin{document}
\maketitle

\begin{abstract}
This memo explains the fraud-detection methodology behind the current
\code{agentic-payments} prototype.  It surveys how fraud detection is
typically done in Web2 payments, ecommerce, and digital banking; identifies
which ideas transfer to agentic payment rails; and specifies the exact
feature, clustering, scoring, explanation, and evaluation methods implemented
in the repository.  The central claim is that agentic payment risk is primarily
a statistical and graph-learning problem, not a large-language-model problem:
agents leave machine-shaped behavioural fingerprints in timing, gas
optimization, method calls, counterparties, and coordination structure.
\end{abstract}

\section{Executive Summary}

Modern fraud systems combine four ingredients: (i) dense event data, (ii)
behavioural features, (iii) supervised or semi-supervised risk models, and
(iv) an operational decision layer for review, step-up verification, blocking,
or liability shift.  Stripe Radar, Shopify Fraud Analysis/Protect, Riskified,
Sift, Forter, Sardine, Monzo, and Revolut differ in surface area, but they all
converge on that architecture: real-time scoring, network effects, rules plus
machine learning, explainability for analysts, and feedback loops from
chargebacks or confirmed fraud.

Agentic payments preserve this structure but change the signal.  Human-centric
features such as browser fingerprints, shipping address mismatch, CVC checks,
email reputation, device biometrics, and human time-of-day patterns become
weaker or unavailable.  New agentic features become available: inter-arrival
regularity, 24-hour uniformity, gas-to-basefee tightness, repeated calldata,
contract/method concentration, late counterparty drift, and wallet-to-wallet
coordination.  Because labels are initially sparse, the first product should
not pretend to be a fully supervised fraud classifier.  It should be a
behavioural risk intelligence layer: rank anomalous wallets, explain the
signals, surface coordinated clusters, and improve as real labels arrive.

The repository currently implements:

\begin{itemize}
  \item A synthetic labelled data generator for humans, deterministic agents,
        payment bots, compromised agents, and collusion rings.
  \item A public x402/Base data collector and decoder for EIP-3009-style USDC
        \code{transferWithAuthorization} settlements.
  \item A 36-feature behavioural fingerprint per payer wallet.
  \item Pairwise wallet clustering using counterparty overlap, method
        agreement, and timing co-movement.
  \item An eight-factor composite score calibrated against a low-volume
        baseline, mapped to risk tiers with human-readable explanations.
  \item A dashboard that presents triage, evidence, behaviour maps, and
        coordination rather than just a scatterplot.
\end{itemize}

\section{What Fraud Detection Looks Like in Web2}

\subsection{Recurring Structure}

The fraud literature has been consistent for two decades: fraud is rare,
adaptive, adversarial, and operationally asymmetric.  Bolton and Hand's classic
review frames statistical fraud detection as the fallback once prevention
fails, and emphasizes that fraudsters adapt to controls over time
\cite{bolton2002}.  Ngai et al. classify financial fraud detection work around
data mining families such as classification, clustering, outlier detection,
regression, visualization, and social-network analysis \cite{ngai2011}.
Later credit-card fraud work stresses feature engineering, class imbalance,
delayed labels, and concept drift \cite{bahnsen2016,dalpozzolo2015,carcillo2018}.

In production systems, the common architecture is:

\begin{enumerate}
  \item \textbf{Event capture:} payment, account, device, browser, network,
        order, chargeback, customer-service, and identity events.
  \item \textbf{Feature computation:} real-time, near-real-time, and batch
        aggregations at transaction, account, device, card, merchant, and
        counterparty levels.
  \item \textbf{Decisioning:} rules, supervised models, anomaly detectors,
        graph/link analysis, and risk thresholds.
  \item \textbf{Action layer:} allow, review, challenge, step-up verification,
        block, reserve funds, cancel order, or reimburse/guarantee merchant.
  \item \textbf{Feedback loop:} chargebacks, dispute outcomes, manual review
        decisions, confirmed scams, law-enforcement alerts, and analyst labels.
\end{enumerate}

\subsection{Stripe Radar}

Stripe Radar is the closest Web2 analogue to our intended product.  Stripe
documents risk evaluation as a machine-learning score and risk level for each
payment.  Radar for Fraud Teams exposes a 0--99 risk score, with default
thresholds around 65 for elevated risk and 75 for high risk \cite{stripeRisk}.
Radar risk insights describe an adaptive AI model that uses hundreds of
signals and data across Stripe's business network, and exposes some explanatory
signals to analysts \cite{stripeInsights}.  Radar rules then let merchants
block or review payments according to score, risk level, CVC/AVS checks, and
custom predicates \cite{stripeRules}.  Radar for Platforms extends the same
logic from transaction fraud to connected-account risk: machine-learning risk
scores, customized rules, analyst investigation, document/selfie verification,
and actions such as reserves \cite{stripePlatforms}.

\textbf{Methodological lesson:} the product is not ``a model.''  It is a
risk operating system: score, explain, rule, review, and act.  Our dashboard
should therefore show queue, reason, action, and feedback hooks, not merely an
embedding plot.

\subsection{Shopify Fraud Analysis and Protect}

Shopify's fraud analysis gives order-level fraud recommendations and indicators
for online card orders, and its documentation states that the recommendations
are powered by machine-learning algorithms trained on historical transactions
across Shopify stores \cite{shopifyFraud}.  Shopify Protect adds a business
model layer: for eligible Shop Pay transactions, Shopify can cover protected
fraudulent/unrecognized chargebacks and manage the dispute process
\cite{shopifyProtect}.

\textbf{Methodological lesson:} ecommerce fraud systems are not just score
providers; they change merchant behaviour by saying which orders are safe to
fulfil, which require verification, and where liability sits.  Our equivalent
actions are: allow, monitor, step-up auth, throttle/sandbox agent, freeze, or
submit dispute/AML evidence.

\subsection{Chargeback-Guarantee and Risk-Network Vendors}

Riskified describes real-time approve/decline decisions from machine-learning
models over hundreds of transaction features, a merchant network with over a
billion historical transactions, and an aligned-incentive chargeback guarantee
\cite{riskifiedGuarantee}.  Its product pages emphasize data collection,
device/browser behavioural insight, merchant-network linking, and explainable
machine learning \cite{riskifiedSolutions}.  Sift advertises a global data
network of roughly one trillion annual events, hundreds of brands, rulesets,
and split-second decisioning \cite{sift}.  Forter emphasizes identity
protection and a global identity graph across more than a billion online
identities \cite{forter}.  Sardine combines consortium data, supervised ML,
device intelligence, behaviour biometrics, graph connections, rules, and case
management \cite{sardineDocs,sardineHome}.

\textbf{Methodological lesson:} the defensible asset is a labelled network of
behaviour, not a one-off classifier.  Our wedge should start with wallet-level
behavioural features and grow into an agent reputation graph.

\subsection{Digital Banks and Real-Time Payment Fraud}

Digital banks make the latency problem explicit.  Monzo's fraud platform
separates feature computation from control execution, uses just-in-time,
near-real-time, and batch features, and records control metadata to measure
effectiveness and false positives \cite{monzoReactive}.  Monzo's ML stack also
emphasizes reusable feature engineering and feature monitoring
\cite{monzoML}.  Revolut's scam-detection feature uses machine learning to
detect likely scam payments, decline high-risk payments, send users through an
intervention flow, ask additional questions, and escalate to fraud specialists
\cite{revolutScam}.

\textbf{Methodological lesson:} risk models should be measured as controls on a
payment hot path.  The key product metrics are not only AUC or accuracy, but
latency, recall at tolerable false-positive rate, false-positive cost,
queue volume, and loss prevented.

\section{Why Agentic Payments Are Different}

Agentic payments are not merely card payments with a new front end.  Protocols
such as x402 make payments HTTP-native and machine-executable.  Coinbase's
documentation describes facilitators that verify signed payment payloads and
settle them on-chain, with hosted support for stablecoin payments and EIP-3009
compatible tokens on Base/Polygon and SPL tokens on Solana \cite{x402Facilitator}.
x402's flow lets a server respond with HTTP 402, the client resubmit with a
payment payload, and a facilitator verify and settle before the server returns
the requested resource \cite{x402How}.  This yields a new machine-to-machine
payment substrate.

The fraud signal changes in four ways:

\begin{enumerate}
  \item \textbf{Human context weakens.} Many Web2 signals -- browser history,
        shipping address, email reputation, human typing patterns, and
        shopper/customer-service history -- are less relevant for autonomous
        agents.
  \item \textbf{Machine regularity strengthens.} Agents are more likely to
        transact at high frequency, on cron-like intervals, across all hours,
        with repeated method selectors and narrow value bands.
  \item \textbf{On-chain settlement exposes execution details.} On Base/x402
        we can see facilitator/signer addresses, USDC contract calls,
        calldata method selectors, gas used, gas price, payer, recipient, and
        settlement time.
  \item \textbf{Risk propagates through agent networks.} Compromised agents
        can call other agents, APIs, wallets, or merchants.  This motivates
        graph and cascade methods rather than purely per-transaction scoring.
\end{enumerate}

Therefore the product should begin as \emph{behavioural risk intelligence for
agent traffic}: wallet fingerprints, anomaly factors, coordination clusters,
and interpretable triage.  Fully supervised fraud labels can arrive later via
chargebacks, confirmed compromised wallets, facilitator declines, AML alerts,
merchant disputes, and manual analyst review.

\section{Data Model}

\subsection{Canonical Transaction Row}

The repository standardizes both synthetic and real x402 data to the following
schema:

\begin{center}
\begin{tabular}{ll}
\toprule
Field & Meaning \\
\midrule
\code{block_time_s} & Seconds from window start \\
\code{block_time} & Timestamp when available \\
\code{from_addr} & Payer wallet, agent wallet, or synthetic user wallet \\
\code{to_addr} & Recipient, merchant, contract, or counterparty \\
\code{value_usd} & Transaction value normalized to USD/USDC units \\
\code{gas_used} & Execution gas used \\
\code{gas_price_gwei} & Gas price converted to gwei \\
\code{method_id} & EVM function selector, e.g. \code{0xe3ee160e} \\
\code{success} & Boolean transaction success flag \\
\bottomrule
\end{tabular}
\end{center}

Synthetic data additionally has ground-truth labels:
\code{human_random}, \code{agent_arb_deterministic}, \code{agent_payment_bot},
\code{agent_compromised}, and \code{collusion_ring}.  Public x402/Base data is
decoded from facilitator/signer transactions to Base USDC.  In the current
pull, we normalize 7,512 facilitator-address transactions, decode 437
x402/EIP-3009 USDC payment rows, and score 116 payer wallets.

\subsection{Why Synthetic Data Is Legitimate Here}

Synthetic data is not evidence that real fraud exists.  It is a controlled
instrumentation test.  We need to know whether the pipeline can recover known
policies:

\begin{itemize}
  \item Humans: sparse, diurnal, heterogeneous values and counterparties.
  \item Deterministic agents: periodic, 24-hour, gas-tight, low method entropy.
  \item Payment bots: bursty but consistent recipients and transfer methods.
  \item Compromised agents: late drift to new counterparties and higher burst
        rate/value changes.
  \item Collusion rings: multiple wallets transacting with the same target in
        tight time windows.
\end{itemize}

The public x402 overlay then answers a separate question: does the schema and
feature machinery work on real settlement data?

\section{Feature Engineering}

For every wallet \(w\), let \(T_w=\{t_1,\ldots,t_n\}\) be sorted transaction
times, \(C_w\) counterparties, \(M_w\) method selectors, \(V_w\) values, and
\(G_w\) gas observations.  We compute 36 features in six families.

\subsection{Temporal Features}

Inter-arrival times are
\[
  \Delta_i = t_i - t_{i-1}, \quad i=2,\ldots,n.
\]
The coefficient of variation is
\[
  \mathrm{CV}_\Delta(w)=
  \frac{\operatorname{std}(\Delta)}{\operatorname{mean}(\Delta)+\epsilon}.
\]
Near-zero values are cron-like; values near one resemble Poisson-like
arrival; large values indicate burst/drift behaviour.  We also compute
Shannon entropy of log-binned inter-arrivals:
\[
  H(\Delta)=-\sum_b p_b \log_2 p_b.
\]
Time-of-day uniformity is measured as a Kolmogorov--Smirnov distance between
the empirical hour-of-day CDF and a uniform 24-hour CDF:
\[
  D_{\mathrm{uniform}}(w)=\sup_h |F_w(h)-F_U(h)|.
\]
Low \(D_{\mathrm{uniform}}\) is suspicious for agents because always-on agents
can cover the full day evenly; humans tend to be diurnal.  We also keep a
KL-divergence from a canonical human time-of-day curve.

Burstiness follows Goh and Barabasi's coefficient:
\[
  B(w)=\frac{\sigma_\Delta-\mu_\Delta}{\sigma_\Delta+\mu_\Delta+\epsilon}.
\]

\subsection{Gas and Optimization Features}

Agents are expected to optimize execution more tightly than humans.  We compute
mean and standard deviation of gas price, gas-to-basefee ratio, gas-used
coefficient of variation, and a simple optimization score:
\[
  \mathrm{gas\_opt}(w)=
  \frac{1}{1+\operatorname{std}(g/\mathrm{basefee})+\mathrm{CV}(\mathrm{gas\_used})}.
\]
In synthetic data basefee is simulated; in real Base data, gas price and gas
used are directly pulled from Blockscout.

\subsection{Value, Counterparty, and Method Features}

Value features include mean, median, standard deviation, log-standard
deviation, entropy, and the share of transactions above the population 95th
percentile.  Counterparty features include number of unique counterparties,
top-1 and top-3 concentration, entropy, repeat rate, and late-new-counterparty
share:
\[
  \mathrm{late\_new}(w)=
  \frac{\#\{ \text{transactions in final quartile to previously unseen } c\}}
       {\#\{\text{transactions}\}}.
\]
This is the primary compromised-agent drift signal.

Method features include unique method count, dominant selector share, native
transfer share, ERC-20 transfer share, and success rate.  Repeated calls to a
single selector, for example USDC \code{transferWithAuthorization}, are normal
for a payment facilitator but can still be a useful concentration feature in
combination with timing and counterparty features.

\subsection{Coordination Features}

We compute the peak count in any 60-second window, burst windows per hour, mean
inter-burst minutes, and the number of distinct partner wallets with at least
one transaction within a 60-second window of the wallet's own transactions.
These features are intentionally cohort-dependent: they identify wallets whose
behaviour is anomalous relative to the other wallets in the same observed
window.

\section{Wallet Coordination Graph}

For each wallet \(w\), define a profile
\[
  P_w = \{(c,m,t): \text{wallet } w \text{ called counterparty } c
          \text{ with method } m \text{ at time } t\}.
\]
For wallet pair \((i,j)\), compute:

\begin{align*}
O_{ij} &= \frac{|C_i\cap C_j|}{|C_i\cup C_j|},\\
M_{ij} &= \frac{1}{|C_i\cap C_j|}\sum_{c\in C_i\cap C_j}
          \mathbf{1}\{M_i(c)\cap M_j(c)\neq \varnothing\},\\
T_{ij} &= \frac{\#\{(t_i,t_j): |t_i-t_j|\leq 60s,\ c_i=c_j\}}
              {\#\{(t_i,t_j): c_i=c_j\}}.
\end{align*}

The composite similarity is
\[
  S_{ij}=0.3O_{ij}+0.4M_{ij}+0.3T_{ij}.
\]
An edge is created when \(O_{ij}\geq 0.10\) and \(S_{ij}\geq 0.40\).  Connected
components with at least two connected wallets form coordination clusters.
The pairwise loop is capped at the top 500 wallets by transaction count to keep
the hackathon implementation tractable.

This mirrors the Polymarket wallet clusterer, replacing market overlap and
direction agreement with counterparty overlap and method agreement.

\section{Composite Risk Score}

\subsection{Baseline Calibration}

Because labels are sparse in real x402 data, we use a low-volume baseline as a
proxy for ordinary human/low-automation behaviour:
\[
  B=\{w: \mathrm{tx\_count}(w)\leq Q_{0.30}(\mathrm{tx\_count})\}.
\]
For feature \(x\), define baseline z-score:
\[
  z_x(w)=\frac{x(w)-\mu_B(x)}{\sigma_B(x)+\epsilon}.
\]

Some factors are suspicious when high; others are suspicious when low.  We map
z-scores to a 0--100 one-sided suspicion score:
\[
  q_{\mathrm{high}}(z)=
  100\cdot \frac{1}{1+\exp(-1.5(z-1))},
  \quad
  q_{\mathrm{low}}(z)=q_{\mathrm{high}}(-z).
\]
The shift by one standard deviation keeps baseline observations below the
medium-risk threshold and makes high scores require genuine tail behaviour.

\subsection{Eight Factors and Weights}

\begin{center}
\begin{tabular}{lll}
\toprule
Factor & Direction & Weight \\
\midrule
Inter-arrival anomaly & large deviation from baseline median & 0.15 \\
Time-of-day uniformity & unusually uniform, low KS & 0.15 \\
Gas tightness & unusually low gas variance & 0.20 \\
Counterparty concentration & high top-1 share with volume & 0.10 \\
Method concentration & high top selector share with volume & 0.05 \\
Burst intensity & high bursts per hour & 0.10 \\
Drift signal & high late-new-counterparty share & 0.10 \\
Coordination & high partner count in short windows & 0.15 \\
\bottomrule
\end{tabular}
\end{center}

The composite score is
\[
  \risk(w)=
  \operatorname{clip}_{[0,100]}\left(
    \sum_{k=1}^{8}\alpha_k q_k(w)
  \right).
\]
Risk tiers are:
\[
  \text{critical}\geq 75,\quad
  \text{high}\geq 50,\quad
  \text{medium}\geq 25,\quad
  \text{low}<25.
\]

This follows the production pattern used by Radar-like tools: scalar score,
risk tier, rule/action layer, and analyst explanation.

\subsection{Explanations}

For each wallet, we sort the eight factor scores and generate a short
explanation from the top three.  Examples:

\begin{itemize}
  \item ``inter-arrival CV 0.05 (cron); 100\% calls to one selector; 36\%
        volume to single counterparty.''
  \item ``inter-arrival CV 1.45 (bursty); 98\% volume to single counterparty;
        100\% calls to one selector.''
\end{itemize}

This deliberately mirrors Web2 risk insights: the model should not simply
declare a wallet risky; it should say what behavioural facts drove the score.

\section{Embedding and Dashboard}

The embedding is not the primary detector.  It is an analyst visualization.
We standardize numeric features and project them to two dimensions with UMAP,
falling back to PCA if UMAP is unavailable.  The dashboard now leads with a
triage queue and selected-case evidence, then offers the embedding as a
behaviour map.  This follows the Polymarket dashboard pattern: overview cards,
ranked tables, badges, tabs, factor bars, and evidence before visualization.

\section{Current Results}

\subsection{Synthetic Stress Test}

The synthetic run has 52,162 transactions and 188 wallet rows.  The current
calibration flags deterministic arbitrage agents and collusion rings as high
risk, catches most compromised agents, and keeps human wallets out of the
high-risk queue.  This validates that the pipeline can recover the behaviours
it is designed to detect.

\subsection{Public x402/Base Overlay}

The no-key public data collector currently pulls:

\begin{itemize}
  \item 24 facilitators from x402.watch.
  \item 43 x402-related addresses/signers.
  \item 7,512 Base facilitator-address transactions.
  \item 437 decoded x402/EIP-3009 USDC payment rows.
  \item 116 payer wallets and 19 recipient wallets.
\end{itemize}

The public x402 payer-wallet scorer currently returns 65 medium-risk wallets,
51 low-risk wallets, and no high-risk claims.  This is an important honest
result: the real public slice is usable, but small and weakly labelled.  It
should be presented as a real-data overlay, not as proof of fraud.

\section{Why These Methods, Not an LLM Classifier?}

LLMs are useful for analyst summaries, documentation, and case narratives.  The
core fraud decision should be statistical and graph-based because:

\begin{itemize}
  \item Fraud decisions require stable calibration and measurable false-positive
        rates.
  \item Payment risk depends on structured features and temporal/graph context.
  \item Regulators, merchants, and banks need explanations grounded in observed
        facts.
  \item New fraud modes produce drift and adversarial adaptation; feature and
        graph monitoring are easier to audit than generative outputs.
\end{itemize}

Future versions may use LLMs to generate analyst-readable case summaries, but
the score should remain feature-grounded.

\section{Evaluation Plan}

The evaluation should mature in four stages:

\begin{enumerate}
  \item \textbf{Synthetic recovery:} policy-level precision/recall for known
        deterministic agents, compromised agents, humans, and collusion rings.
  \item \textbf{Real overlay sanity checks:} distribution of x402 features,
        duplicate/failed settlements, top recipients, facilitator drift, and
        stability of score thresholds over time.
  \item \textbf{Weak labels:} Forta alerts, sanctioned/high-risk address lists,
        facilitator KYT declines, merchant disputes, and manually verified
        suspicious wallets.
  \item \textbf{Supervised calibration:} cost-sensitive models, precision-recall
        curves, temporal holdouts, delayed-label handling, and active learning
        for analyst review.
\end{enumerate}

Metrics should include precision at queue size, recall at fixed false-positive
rate, expected loss avoided, review burden, latency, and threshold stability.
For real-time deployment, the relevant comparison is not only model AUC but
whether a control can safely run on the payment path.

\section{Future Model Roadmap}

\begin{enumerate}
  \item \textbf{Feature store:} separate just-in-time features from
        near-real-time and batch aggregates, following the pattern used in
        digital-bank fraud systems.
  \item \textbf{Graph model:} build a wallet--recipient--facilitator graph and
        use graph embeddings or GNNs once labels exist.  CARE-GNN and related
        camouflage-resistant methods are relevant because fraudsters adapt both
        attributes and edges \cite{caregnn}.  Blockchain AML work on Elliptic
        datasets demonstrates the utility of temporal transaction graphs
        \cite{ellipticGcn,elliptic2}.
  \item \textbf{Reputation layer:} convert wallet risk into portable agent,
        facilitator, merchant, and counterparty reputation signals.
  \item \textbf{Decision policies:} map tiers to concrete actions: allow,
        monitor, step-up authorization, rate-limit, sandbox agent, freeze,
        request identity/delegation proof, or create dispute/AML evidence.
  \item \textbf{Feedback loop:} store analyst decisions and post-event labels
        so the model compounds data advantage over time.
\end{enumerate}

\section{Generalising to Traditional Payment Rails}

The current x402/Base implementation should be treated as the first rail
adapter, not as the boundary of the product.  The invariant is not a wallet; it
is an actor producing payment behaviour.  A production architecture should
therefore separate rail-specific ingestion from a common event/entity graph.

\subsection{Rail Adapters}

Stripe is the easiest traditional-payment adapter because the Charge object
already exposes risk outcomes, network status, decline/advice codes,
\code{fraud_details}, card checks, payment-method fingerprints, refunds,
disputes, and payment-method-specific details \cite{stripeCharge}.  Stripe's
Radar score and risk level should be stored as external model signals, not as
ground truth.

Shopify is the natural ecommerce adapter.  Its order-level fraud analysis
surfaces AVS/CVV, IP, multiple-card attempts, and low/medium/high
recommendations; the Order Risk API allows third-party risk assessments to be
displayed back to merchants with accept/investigate/cancel recommendations
\cite{shopifyFraud,shopifyOrderRisk}.  This makes Shopify a practical path to
a private app demo: ingest orders and transactions, produce an agent-aware
order-risk record, and write the explanation into the merchant workflow.

Visa and Mastercard direct integrations are not hackathon-easy because their
risk-management products are partner/restricted surfaces.  Public Visa material
describes Visa Advanced Authorization and Visa Risk Manager as in-flight AI
risk scoring plus rules/case-management tooling over VisaNet-scale attributes
\cite{visaVaa}.  Mastercard describes Decision Intelligence and Transaction
Fraud Monitoring as real-time/pre-authorization risk scoring products with
rules, case management, business insights, and network intelligence
\cite{mastercardDi,mastercardTfm}.  The practical route is through an issuer,
acquirer, PSP, or payment facilitator that can provide authorization streams
and chargeback labels.

\subsection{Canonical Entity Graph}

The wallet-centric schema should generalise to:

\begin{itemize}
  \item \textbf{Payment events:} rail, time, amount, currency, status, payer,
        payee, instrument, merchant, order, channel, external score, and raw
        rail reference.
  \item \textbf{Entities:} customer/account, wallet, card token/fingerprint,
        bank account fingerprint, merchant, device/session, email/phone,
        IP/ASN, shipping/billing address hash, and agent identity where
        available.
  \item \textbf{Labels:} chargeback, confirmed fraud, issuer decline, Radar
        high/elevated/normal, Shopify high/medium/low, refund abuse, scam
        report, AML alert, and manual-review decision.
\end{itemize}

\subsection{Feature Transfer}

Most x402 features have direct traditional-payment analogues:

\begin{center}
\begin{tabular}{ll}
\toprule
x402 feature family & Traditional-payment equivalent \\
\midrule
Inter-arrival regularity & checkout/payment-attempt velocity and retry cadence \\
24-hour uniformity & always-on API or checkout automation \\
Gas tightness & repeated request/device/API-idempotency patterns \\
Counterparty concentration & merchant, SKU, recipient, or address concentration \\
Method concentration & card/wallet/BNPL/3DS strategy concentration \\
Burst intensity & card testing, checkout waves, promo/refund abuse \\
Late counterparty drift & ATO or compromised-agent shift to new merchants/SKUs \\
Coordination & linked accounts/cards/devices/IPs in tight windows \\
\bottomrule
\end{tabular}
\end{center}

Traditional rails also add features that x402 lacks: AVS/CVV outcomes, 3DS/ECI
signals, device/browser entropy, shipping/billing distance, merchant-category
risk, issuer/network decline codes, capture delay, refund velocity,
fulfilment status, and chargeback reason codes.

The venture-level framing is therefore: \emph{behavioural risk intelligence for
automated payment traffic across rails}.  We start with x402 because it makes
machine payments visible, then reuse the same fingerprinting and graph-risk
engine for Stripe merchants, Shopify stores, PSPs, acquirers, issuers, and
digital banks.

\section{Limitations}

\begin{itemize}
  \item The synthetic labels are useful for pipeline validation but not proof of
        real-world fraud.
  \item The public x402 data pull is small and weakly labelled; no high-risk
        claim should be made from it yet.
  \item The low-volume baseline is a pragmatic proxy, not a guaranteed human
        baseline.
  \item Gas features need chain-specific normalization when moving beyond Base.
  \item Counterparty concentration can be benign for legitimate API/payment
        agents; it is only suspicious in combination with timing, drift, gas,
        and coordination.
  \item A production system needs delayed-label handling, threshold governance,
        monitoring, and analyst feedback.
\end{itemize}

\section{Implementation Map}

\begin{center}
\begin{tabular}{ll}
\toprule
Method & Repository path \\
\midrule
Synthetic policy generator & \code{src/ingest/synthetic.py} \\
Public x402 collector/decoder & \code{src/ingest/public_x402.py} \\
Behavioural features & \code{src/features/fingerprint.py} \\
Wallet clustering & \code{src/models/cluster.py} \\
Composite scoring & \code{src/models/score.py} \\
Embedding & \code{src/viz/embed.py} \\
Demo dashboard & \code{src/viz/dashboard.py} \\
Synthetic pipeline & \code{src/pipeline.py} \\
Public x402 pipeline & \code{src/pipeline_public.py} \\
\bottomrule
\end{tabular}
\end{center}

\begin{thebibliography}{99}

\bibitem{bolton2002}
Richard J. Bolton and David J. Hand.
\newblock Statistical Fraud Detection: A Review.
\newblock \emph{Statistical Science}, 17(3):235--255, 2002.
\newblock \url{https://projecteuclid.org/journals/statistical-science/volume-17/issue-3/Statistical-Fraud-Detection-A-Review/10.1214/ss/1042727940.pdf}

\bibitem{ngai2011}
E. W. T. Ngai, Yong Hu, Y. H. Wong, Yijun Chen, and Xin Sun.
\newblock The application of data mining techniques in financial fraud detection: A classification framework and an academic review of literature.
\newblock \emph{Decision Support Systems}, 50(3):559--569, 2011.
\newblock \url{https://doi.org/10.1016/j.dss.2010.08.006}

\bibitem{bahnsen2016}
Alejandro Correa Bahnsen, Djamila Aouada, Aleksandar Stojanovic, and Bjoern Ottersten.
\newblock Feature Engineering Strategies for Credit Card Fraud Detection.
\newblock \emph{Expert Systems with Applications}, 51:134--142, 2016.
\newblock \url{https://doi.org/10.1016/j.eswa.2015.12.030}

\bibitem{dalpozzolo2015}
Andrea Dal Pozzolo, Olivier Caelen, Reid A. Johnson, and Gianluca Bontempi.
\newblock Calibrating Probability with Undersampling for Unbalanced Classification.
\newblock IEEE Symposium Series on Computational Intelligence, 2015.
\newblock \url{https://doi.org/10.1109/SSCI.2015.33}

\bibitem{carcillo2018}
Fabrizio Carcillo, Andrea Dal Pozzolo, Yann-Ael Le Borgne, Olivier Caelen,
Yannis Mazzer, and Gianluca Bontempi.
\newblock Streaming active learning strategies for real-life credit card fraud detection.
\newblock \emph{International Journal of Data Science and Analytics}, 5:285--300, 2018.
\newblock \url{https://link.springer.com/article/10.1007/s41060-018-0116-z}

\bibitem{stripeRisk}
Stripe Documentation.
\newblock Risk evaluation: risk score and risk outcomes.
\newblock \url{https://docs.stripe.com/radar/risk-evaluation}

\bibitem{stripeCharge}
Stripe Documentation.
\newblock The Charge object.
\newblock \url{https://docs.stripe.com/api/charges/object}

\bibitem{stripeInsights}
Stripe Documentation.
\newblock Risk insights.
\newblock \url{https://docs.stripe.com/radar/reviews/risk-insights}

\bibitem{stripeRules}
Stripe Documentation.
\newblock Fraud prevention rules.
\newblock \url{https://docs.stripe.com/radar/rules}

\bibitem{stripePlatforms}
Stripe Documentation.
\newblock Radar for Platforms.
\newblock \url{https://docs.stripe.com/connect/risk-management/risk-tooling}

\bibitem{shopifyFraud}
Shopify Help Center.
\newblock Fraud analysis.
\newblock \url{https://help.shopify.com/en/manual/fulfillment/managing-orders/protecting-orders/fraud-analysis}

\bibitem{shopifyOrderRisk}
Shopify Developer Documentation.
\newblock Order Risk resource.
\newblock \url{https://shopify.dev/docs/api/admin-rest/2025-07/resources/order-risk}

\bibitem{shopifyProtect}
Shopify.
\newblock Shopify Protect.
\newblock \url{https://www.shopify.com/protect}

\bibitem{visaVaa}
Visa.
\newblock Visa Advanced Authorization + Visa Risk Manager: Managing fraud better, together.
\newblock \url{https://corporate.visa.com/content/dam/VCOM/corporate/solutions/documents/visa-risk-manager-cnp-global-infographic.pdf}

\bibitem{mastercardDi}
Mastercard.
\newblock Decision Intelligence for fraud and risk management.
\newblock \url{https://www.mastercard.com/global/en/business/cybersecurity-fraud-prevention/risk-decisioning/decision-intelligence.html}

\bibitem{mastercardTfm}
Mastercard.
\newblock Transaction Fraud Monitoring.
\newblock \url{https://www.mastercard.com/global/en/business/cybersecurity-fraud-prevention/risk-decisioning/transaction-fraud-monitoring.html}

\bibitem{riskifiedGuarantee}
Riskified.
\newblock Chargeback Guarantee and Fraud Protection.
\newblock \url{https://www.riskified.com/chargeback-guarantee/}

\bibitem{riskifiedSolutions}
Riskified.
\newblock Ecommerce fraud prevention solutions.
\newblock \url{https://www.riskified.com/fraud/prevention-solutions/}

\bibitem{sift}
Sift.
\newblock Digital Fraud Prevention and Risk-Based Authentication.
\newblock \url{https://sift.com/}

\bibitem{forter}
Forter.
\newblock Account Protection.
\newblock \url{https://www.forter.com/platform/identity-protection/}

\bibitem{sardineDocs}
Sardine Documentation.
\newblock What Powers Sardine.
\newblock \url{https://docs.sardine.ai/guides/public/gettingstarted/whatpowerssardine}

\bibitem{sardineHome}
Sardine.
\newblock Agentic Financial Crime Platform for Fraud Prevention and AML.
\newblock \url{https://www.sardine.ai/}

\bibitem{monzoReactive}
Monzo.
\newblock Building a reactive Fraud Prevention Platform.
\newblock \url{https://monzo.com/blog/build-a-reactive-fraud-prevention-platform}

\bibitem{monzoML}
Monzo.
\newblock Monzo's machine learning stack.
\newblock \url{https://monzo.com/blog/2022/04/26/monzos-machine-learning-stack}

\bibitem{revolutScam}
Revolut.
\newblock Revolut launches AI feature to protect customers from card scams.
\newblock \url{https://www.revolut.com/en-US/news/revolut_launches_ai_feature_to_protect_customers_from_card_scams_and_break_the_scammers_spell/}

\bibitem{x402Facilitator}
Coinbase Developer Platform.
\newblock x402 Facilitator.
\newblock \url{https://docs.cdp.coinbase.com/x402/core-concepts/facilitator}

\bibitem{x402How}
Coinbase Developer Platform.
\newblock How x402 Works.
\newblock \url{https://docs.cdp.coinbase.com/x402/core-concepts/how-it-works}

\bibitem{caregnn}
Yingtong Dou, Zhiwei Liu, Li Sun, Yuening Deng, Hao Peng, and Philip S. Yu.
\newblock Enhancing Graph Neural Network-based Fraud Detectors against Camouflaged Fraudsters.
\newblock CIKM, 2020.
\newblock \url{https://arxiv.org/abs/2008.08692}

\bibitem{ellipticGcn}
Mark Weber et al.
\newblock Anti-Money Laundering in Bitcoin: Experimenting with Graph Convolutional Networks for Financial Forensics.
\newblock KDD Workshop on Anomaly Detection in Finance, 2019.
\newblock \url{https://arxiv.org/abs/1908.02591}

\bibitem{elliptic2}
Elliptic, MIT, and IBM researchers.
\newblock The Shape of Money Laundering: Subgraph Representation Learning on the Blockchain with the Elliptic2 Dataset.
\newblock 2024.
\newblock \url{https://arxiv.org/abs/2404.19109}

\end{thebibliography}

\end{document}
